%!TEX root = ../main.tex
% Resumen en español. Las palabras clave se definen en main.tex.

\begin{resumen}

El cuidado de una persona con discapacidad suele recaer en familiares sin
formación específica, en cuidadores contratados o en el personal de una
institución. Quien cuida tiene que organizar medicación, turnos, estudios y
trámites, y la información que necesita está repartida entre papeles,
mensajes y sitios oficiales difíciles de recorrer. Cuando cuidan varias
personas, además, no existe un registro común de qué se hizo y quién lo
hizo.

Este trabajo presenta Acompañar, una plataforma web adaptada al celular que
centraliza la gestión del cuidado y organiza a quienes participan de él en
grupos con un dueño, cuidadores que registran el día a día y observadores
que solo consultan. El componente central es un sistema de recomendación
híbrido y explicable: considera solo los recursos del tipo de discapacidad
de la persona, compara su perfil con un catálogo verificado, descarta lo que
no corresponde, incorpora las valoraciones de otros cuidadores y explica
cada sugerencia. Se suman una guía de trámites del Certificado Único de
Discapacidad basada en normativa oficial, un módulo de emergencia y un modelo
de suscripción con un plan gratuito y dos planes pagos. El sistema se
construyó con React, ASP.NET Core y SQL Server; hoy se ejecuta en un entorno
local y su publicación en Azure queda planteada como próximo paso.

La evaluación incluyó 68 pruebas automatizadas, el análisis de las
recomendaciones sobre trece perfiles de demostración, con el 100~\% de las
sugerencias explicadas, la verificación de 52 enlaces contra fuentes
oficiales y un relevamiento analizado sobre un estudio simulado de 80 casos.
El aporte del trabajo es la integración, en un solo lugar y adaptada al
contexto argentino, de la organización del cuidado, la coordinación del
grupo y una orientación explicable. La validación con usuarios reales queda
como el paso siguiente.

\end{resumen}
